\chapter{Probability}
\label{ch:probability}

Econometric estimators are functions of random samples.  Probability gives
us the language for describing their behavior before the sample is observed.
This chapter reviews the ideas used repeatedly later in the course.  Optional
measure-theoretic details, longer simulations, and the delta method appear in
Appendix~\ref{app:probability}.

\section{Random variables and distributions}

A random variable is a numerical function of an underlying random outcome.
For example, the outcome might describe all features of tomorrow's weather,
while $X$ records only whether it rains.  The set of values that $X$ can take
with positive probability, or can approach arbitrarily closely, is its
\emph{support}.  The support should not be confused with the underlying
sample space of outcomes.

A random variable is \emph{discrete} when its distribution is concentrated
on a finite or countable set.  Its probability mass function (p.m.f.) is
\[
  p_X(x)=\Pr(X=x).
\]
A Bernoulli random variable $X\sim\operatorname{Bernoulli}(p)$ has support
$\{0,1\}$ and $\Pr(X=1)=p$.

\begin{rcode}
set.seed(5280)
rbinom(n = 10, size = 1, prob = 0.5)
\end{rcode}

The cumulative distribution function (c.d.f.) is defined for every random
variable:
\[
  F_X(x)=\Pr(X\leq x),\qquad x\in\mathbb R.
\]
If $X$ has a density $f_X$, then
\[
  F_X(x)=\int_{-\infty}^{x}f_X(s)\,ds,
  \qquad
  \Pr(a<X\leq b)=\int_a^b f_X(s)\,ds.
\]
Such a variable is called \emph{absolutely continuous}, and
$\Pr(X=x)=0$ for every $x$.  Probability is not generally the geometric
length of an interval: it is the area under the density.  Length gives
probability only for a suitably normalized uniform distribution.

A normal variable $X\sim N(\mu,\sigma^2)$ has density
\[
  f_X(x)=\frac{1}{\sigma\sqrt{2\pi}}
  \exp\!\left\{-\frac{(x-\mu)^2}{2\sigma^2}\right\}.
\]
Normal distributions are important for large-sample approximations; we will
not assume that economic variables themselves are normally distributed.

\subsection{Joint, marginal, and conditional distributions}

The joint distribution of $(X,Y)$ describes their behavior together.  Its
marginal distributions describe $X$ and $Y$ separately.  Conditional
distributions describe one variable within a subpopulation indexed by the
other.

When a joint density exists,
\[
  f_X(x)=\int f_{X,Y}(x,y)\,dy,
  \qquad
  f_{Y\mid X}(y\mid x)=\frac{f_{X,Y}(x,y)}{f_X(x)}
\]
for values of $x$ with $f_X(x)>0$.  Hence
\[
  f_{X,Y}(x,y)=f_{Y\mid X}(y\mid x)f_X(x)
              =f_{X\mid Y}(x\mid y)f_Y(y).
\]
The joint distribution contains the marginals, but the two marginals alone
usually do not reveal how $X$ and $Y$ move together.

\begin{definition}[Independence]
$X$ and $Y$ are independent, written $X\perp Y$, if
\[
  F_{X,Y}(x,y)=F_X(x)F_Y(y)
\]
for every $(x,y)$.  When densities exist, this is equivalent to
$f_{X,Y}(x,y)=f_X(x)f_Y(y)$ almost everywhere.
\end{definition}

If $X\perp Y$, then $g(X)\perp h(Y)$ for measurable functions $g$ and $h$.

\section{Expectation and dispersion}

For a discrete variable, replace the integrals below by sums.  Provided the
relevant integral exists, the expectation is
\[
  \mathbb E(X)=\int x f_X(x)\,dx=\mu_X.
\]
Expectation is linear:
\[
  \mathbb E(aX+bY)=a\mathbb E(X)+b\mathbb E(Y),
\]
whether or not $X$ and $Y$ are independent.

Variance and covariance are
\begin{align*}
  \operatorname{Var}(X)
    &=\mathbb E[(X-\mu_X)^2]
      =\mathbb E(X^2)-\mu_X^2,\\
  \operatorname{Cov}(X,Y)
    &=\mathbb E[(X-\mu_X)(Y-\mu_Y)]
      =\mathbb E(XY)-\mu_X\mu_Y.
\end{align*}
Useful rules include
\begin{align*}
  \operatorname{Var}(aX+bY)
  &=a^2\operatorname{Var}(X)+b^2\operatorname{Var}(Y)
    +2ab\operatorname{Cov}(X,Y),\\
  \operatorname{Cov}(aX+bY,Z)
  &=a\operatorname{Cov}(X,Z)+b\operatorname{Cov}(Y,Z).
\end{align*}
The correlation is
\[
  \rho_{XY}=\frac{\operatorname{Cov}(X,Y)}
  {\sqrt{\operatorname{Var}(X)\operatorname{Var}(Y)}}\in[-1,1]
\]
when both variances are positive.  Correlation is unit-free, but zero
correlation need not imply independence.

For a $k\times1$ random vector $X$, its mean is the vector
$\mu_X=\mathbb E(X)$ and its covariance matrix is
\[
  \Sigma_X
  =\operatorname{Var}(X)
  =\mathbb E[(X-\mu_X)(X-\mu_X)'].
\]
It is symmetric and p.s.d. because, for every $a\in\mathbb R^k$,
\[
  a'\Sigma_Xa=\operatorname{Var}(a'X)\geq0.
\]

\section{Conditional expectation}

Conditional expectation is among the most important objects in this course.
For a continuous pair with a conditional density,
\[
  \mathbb E(Y\mid X=x)=\int y f_{Y\mid X}(y\mid x)\,dy.
\]
For a fixed $x$, this is a number.  By contrast,
$\mathbb E(Y\mid X)$ is a random variable---a function of the still-random
$X$.

Conditional expectation is linear, and variables measurable with respect to
the conditioning information can be taken outside:
\begin{align*}
  \mathbb E(aY+bZ\mid X)
    &=a\mathbb E(Y\mid X)+b\mathbb E(Z\mid X),\\
  \mathbb E[g(X)Y\mid X]
    &=g(X)\mathbb E(Y\mid X).
\end{align*}

\begin{theorem}[Law of iterated expectations]
If the expectations exist, then
\[
  \mathbb E\{\mathbb E(Y\mid X)\}=\mathbb E(Y).
\]
More generally, if the information in $Z$ is contained in the information in
$X$, then
\[
  \mathbb E\{\mathbb E(Y\mid X)\mid Z\}=\mathbb E(Y\mid Z).
\]
\end{theorem}

In particular, $\mathbb E(X\mid X)=X$.  We say that $Y$ is \emph{mean
independent} of $X$ if $\mathbb E(Y\mid X)=\mathbb E(Y)$ almost surely.  Mean
independence implies zero covariance (when second moments exist), but it is
not symmetric in general.  Full independence implies mean independence in
both directions:
\[
  X\perp Y
  \ \Longrightarrow\ 
  \mathbb E(Y\mid X)=\mathbb E(Y)
  \ \Longrightarrow\ 
  \operatorname{Cov}(X,Y)=0.
\]

\section{Random samples and large-sample theory}

For cross-sectional applications, we often assume that
$W_i=(Y_i,X_i')'$ is independent and identically distributed across units
$i$.  This means $W_i\perp W_j$ for $i\neq j$ and each unit has the same
distribution.  It does \emph{not} require the components within $W_i$ to be
independent.  In panel data, repeated observations on the same unit are
normally dependent; i.i.d. sampling may instead be imposed across units or
replaced by a suitable dependence condition.

\subsection{Convergence in probability and the WLLN}

\begin{definition}[Convergence in probability]
A sequence of random vectors $W_n$ converges in probability to $w$, written
$W_n\stackrel{p}{\longrightarrow}w$, if, for every $\varepsilon>0$,
\[
  \Pr(\lVert W_n-w\rVert>\varepsilon)\longrightarrow0.
\]
When $W_n$ is an estimator and $w$ is the target parameter, this property is
called consistency.
\end{definition}

\begin{theorem}[Weak law of large numbers]
Let $X_1,\ldots,X_n$ be i.i.d. random vectors with
$\mathbb E\lVert X_i\rVert<\infty$.  Then
\[
  \bar X_n=\frac1n\sum_{i=1}^n X_i
  \stackrel{p}{\longrightarrow}\mathbb E(X_i).
\]
\end{theorem}

For a scalar random variable with finite variance, the intuition is especially
transparent because
\[
  \operatorname{Var}(\bar X_n)=\frac{\operatorname{Var}(X_i)}{n}\longrightarrow0.
\]

\subsection{Convergence in distribution and the CLT}

\begin{definition}[Convergence in distribution]
$W_n$ converges in distribution to $W$, written
$W_n\stackrel{d}{\longrightarrow}W$, if
\[
  \mathbb E[g(W_n)]\longrightarrow\mathbb E[g(W)]
\]
for every bounded continuous function $g$.  For scalar variables this is
equivalent to $F_{W_n}(w)\to F_W(w)$ at every continuity point of $F_W$.
\end{definition}

It is not correct to require convergence of probabilities for \emph{every}
set: boundary sets can fail even when convergence in distribution holds.

\begin{theorem}[Multivariate central limit theorem]
Let $X_1,\ldots,X_n$ be i.i.d. $k\times1$ random vectors with mean $\mu$ and
finite covariance matrix $\Sigma$.  Then
\[
  \sqrt n(\bar X_n-\mu)
  \stackrel{d}{\longrightarrow}N(0,\Sigma).
\]
\end{theorem}

The limiting distribution is standard normal only after appropriate
standardization.  In the multivariate result, $\Sigma$ contains all
cross-component covariances, not merely the marginal variances.

\subsection{Continuous mapping and Slutsky's theorem}

If $g$ is continuous at the relevant limiting values, the continuous mapping
theorem gives
\[
  W_n\stackrel{p}{\longrightarrow}W
  \Longrightarrow
  g(W_n)\stackrel{p}{\longrightarrow}g(W),
\]
and an analogous statement holds for convergence in distribution.

Slutsky's theorem says that if $X_n\stackrel{d}{\longrightarrow}X$ and
$Y_n\stackrel{p}{\longrightarrow}c$ for a constant $c$, then
\begin{align*}
  X_n+Y_n&\stackrel{d}{\longrightarrow}X+c,\\
  X_nY_n&\stackrel{d}{\longrightarrow}cX,\\
  X_n/Y_n&\stackrel{d}{\longrightarrow}X/c\quad(c\neq0).
\end{align*}
These tools turn probability limits and central limit theorems into the
large-sample distributions of econometric estimators.
